
%----------------------------------------------------------------------------------------
%	PACKAGES AND OTHER DOCUMENT CONFIGURATIONS
%----------------------------------------------------------------------------------------

\documentclass[
	%a4paper, % Uncomment for A4 paper size (default is US letter)
	11pt, % Default font size, can use 10pt, 11pt or 12pt
]{resume} % Use the resume class

%\usepackage{ebgaramond} % Use the EB Garamond font
\usepackage[pdfpagelabels,plainpages=false]{hyperref}
\usepackage{url}
\usepackage{booktabs}
\usepackage{changepage}
\usepackage{enumitem}

\usepackage{fancyhdr}
\pagestyle{fancy}
%\lhead{This is my name}
%\rhead{this is page \thepage}
\cfoot{\thepage}
%\fancyfoot[R]{\fontsize{9}{11}\selectfont \today}
\renewcommand{\headrulewidth}{0.0pt}
%\renewcommand{\footrulewidth}{0.4pt}

\setlist[itemize]{leftmargin=1.3em}
\setlist[enumerate]{leftmargin=1.5em}

%------------------------------------------------

\name{Maxwell Cole, PhD} % Your name to appear at the top

% You can use the \address command up to 3 times for 3 different addresses or pieces of contact information
% Any new lines (\\) you use in the \address commands will be converted to symbols, so each address will appear as a single line.

\address{Medical Physics Resident \textbar\ UC San Diego}
\address{(214)~$\cdot$~934~$\cdot$~1351 \\ maxsicole@gmail.com} % Contact information
\address{maxsicole.github.io}

%----------------------------------------------------------------------------------------

\begin{document}

%----------------------------------------------------------------------------------------
%	RESEARCH AND CLINICAL INTERESTS SECTION
%----------------------------------------------------------------------------------------

\begin{rSection}{Research and Clinical Interests}

My clinical interests are in developing treatment techniques and workflows that can be implemented with equipment and resources clinics already have, including streamlined approaches and automation of manual planning steps. My research interests are in computational and mathematical modeling of radiobiological and biophysical systems, an interest that began during my PhD studying vascular injury and whole-body blood flow at large scale. I am also interested in global health and scientific outreach, building open-access physics education programs and contributing to organizations focused on expanding access to medical care and physics training.

\end{rSection}
\vspace{0.6cm}

%----------------------------------------------------------------------------------------
%	EDUCATION SECTION
%----------------------------------------------------------------------------------------

\begin{rSection}{Education}

	\textbf{Medical Physics Residency} \hfill \textit{Expected Completion:} June 2027 \\
	Department of Radiation Medicine and Applied Sciences \\
	University of California, San Diego

	\textbf{Ph.D. in Physics (Medical Physics)} \hfill 2025 \\
	Louisiana State University
	\vspace{-.2cm}
	\begin{itemize}[label=$\diamond$]
		\itemsep-0.5em
		\item Dissertation: ``First Digital Biophysical Model of the Entire Human Cardiovascular System''
		\item Advisor: Wayne Newhauser, PhD
	\end{itemize}

	\textbf{B.S. in Physics, Minor in Mathematics} \hfill 2020 \\
	Louisiana State University

\end{rSection}
\vspace{0.5cm}

%----------------------------------------------------------------------------------------
%	CERTIFICATIONS SECTION
%----------------------------------------------------------------------------------------

\begin{rSection}{Professional Certifications}

	\textbf{American Board of Radiology (ABR)} \\
	Part 1 -- Completed August 2026 (\textit{Results Pending}); Part 2 -- Scheduled Summer 2027

\end{rSection}
\vspace{0.5cm}

%----------------------------------------------------------------------------------------
%	CLINICAL TRAINING SECTION
%----------------------------------------------------------------------------------------

\begin{rSection}{Clinical Training}

	\textbf{Chief Resident} \hfill 2026-Present
	\vspace{-.2cm}
	\begin{itemize}[label=$\cdot$]
		\itemsep-0.5em
		\item Coordinated call schedules and task assignments for a cohort of 3 co-residents.
		\item Co-authored 3 departmental standard operating procedures.
		\item Served as resident representative to the Quality Committee, presenting safety incidents and developing corrective action strategies to prevent recurrence.
		\item Managed the incident-learning system and throughput tracking across the main site and 4 affiliate clinics.
		\item Served on the Communications Committee, managing the residency program's Instagram account, staffing the residency recruitment booth at the AAPM Annual Meeting, and leading tours and information sessions for prospective applicants.
		\item Trained and mentored junior residents on clinical procedures, quality assurance workflows, and case planning.
	\end{itemize}

\newpage
	\textbf{Treatment Planning and Delivery}
	\vspace{-.2cm}
	\begin{itemize}[label=$\cdot$]
		\itemsep-0.5em
		\item Planned external beam radiotherapy cases spanning all major disease sites using the Eclipse treatment planning system (TPS).
		\item Planned and executed real-time replanning for online adaptive radiotherapy treatments using Ethos.
		\item Planned stereotactic radiosurgery/radiotherapy (SRS/SRT) treatments using both HyperArc and cone-based techniques.
		\item Planned single-field and multi-field optimized (SFO/MFO) proton therapy treatments.
		\item Planned MR-only, synthetic-CT-based prostate VMAT treatments.
		\item Utilized knowledge-based and automated planning tools to improve plan quality and efficiency.
	\end{itemize}

	\textbf{Quality Assurance}
	\vspace{-.2cm}
	\begin{itemize}[label=$\cdot$]
		\itemsep-0.5em
		\item Performed TG-51 dose calibration on Varian TrueBeam, Ethos, and Trilogy linear accelerators.
		\item Performed routine and ad hoc machine QA on Varian TrueBeam, Ethos/Halcyon, and Trilogy linear accelerators; GE CT simulators; Varian Bravos HDR afterloader; MRI; and the Varian ProBeam proton gantry.
		\item Performed pre-treatment stereotactic radiosurgery quality assurance.
		\item Performed patient-specific quality assurance using portal dosimetry, log-file-based verification, and phantom-based measurements.
	\end{itemize}
	
	\textbf{Commissioning and Accreditation}
	\vspace{-.2cm}
	\begin{itemize}[label=$\cdot$]
		\itemsep-0.5em
		\item Participated in commissioning of the AcurosXB dose calculation algorithm (v18) and the ARIA oncology information system upgrade (v18) \textit{(see Clinical Projects)}.
	\end{itemize}

	\textbf{Clinical Responsibilities}
	\vspace{-.2cm}
	\begin{itemize}[label=$\cdot$]
		\itemsep-0.5em
		\item In-vivo dosimetry using MOSFET detectors.
		\item Electron insert factor measurements for custom-fabricated cutouts.
		\item One-on-one physics consults with patients and families on treatment plans and technical questions.
		\item EQD2 conversions for re-treatment cases.
		\item 4DCT simulation and motion management assessment for thoracic and abdominal cases.
		\item Secondary plan checks and weekly chart checks.
	\end{itemize}

	

\end{rSection}
\vspace{0.5cm}

%----------------------------------------------------------------------------------------
%	CLINICAL PROJECTS SECTION
%----------------------------------------------------------------------------------------

\begin{rSection}{Clinical Projects}

	\textbf{Development and Analysis of a Virtual Cone Technique for Trigeminal Neuralgia Treatment with the Millennium MLC} \\
	Supervisor: Ryan Manger, PhD
	\vspace{-.2cm}
	\begin{itemize}[label=$\cdot$]
		\itemsep-0.5em
		\item Optimized and generated a standardized MLC-based planning technique to reproduce the dose distribution of cone-based radiosurgery plans.
		\item Performed film-based dosimetric measurements to validate plan sphericity and conformity against cone-based benchmarks.
		\item Assessed the clinical feasibility of the technique against trigeminal neuralgia dose-specification criteria, identifying dynamic collimator rotation and reduced leaf-gap settings as the primary levers for improving dose sphericity.
		\item \textit{Manuscript and conference presentation planned.}
	\end{itemize}

\newpage
	\textbf{Implementation of a Remote Viewer System for Physician-Covered Special Treatments} \\
	Supervisor: Ryan Manger, PhD
	\vspace{-.2cm}
	\begin{itemize}[label=$\cdot$]
		\itemsep-0.5em
		\item Designed and built a Raspberry Pi-based device to mirror the treatment delivery system console and stream the session to a locally hosted address.
		\item Enabled covering physicians to remotely view on-treatment imaging, contours, and plans without being physically present at the machine console.
	\end{itemize}

	\textbf{Development and Implementation of an Automated Electron Block Procedure} \\
	Supervisor: Ryan Manger, PhD
	\vspace{-.2cm}
	\begin{itemize}[label=$\cdot$]
		\itemsep-0.5em
		\item Developed a program to automate treatment planning for custom electron cutout blocks.
		\item Reduced potential for error by eliminating manual insert-factor measurements and hand calculation of monitor units from the clinical workflow.
	\end{itemize}

	\textbf{Additional Clinical \& Technical Projects}
	\vspace{-.2cm}
	\begin{itemize}[label=$\cdot$]
		\itemsep-0.5em
		\item \textbf{MR-Only Simulation and Planning Workflow} -- Assisted with faculty education and clinical implementation \textit{(Supervisors: T. Pawlicki, PhD; R. Manger, PhD)}
		\item \textbf{AcurosXB v18 Dose Calculation Algorithm} -- Assisted with commissioning and clinical implementation across Ethos, Halcyon, and TrueBeam \textit{(Supervisor: R. Manger, PhD)}
		\item \textbf{ARIA v18 Upgrade} -- Assisted with commissioning and clinical implementation \textit{(Supervisor: R. Manger, PhD)}
		\item \textbf{Annual QA Program, Ethos and Halcyon} -- Designed the annual test and measurement schedule for the clinic's O-ring linear accelerators \textit{(Supervisors: L. Padilla, PhD; J. Hosiak, PhD)}
	\end{itemize}

\end{rSection}
\vspace{0.5cm}

%----------------------------------------------------------------------------------------
%	RESEARCH EXPERIENCE SECTION
%----------------------------------------------------------------------------------------

\begin{rSection}{Research Experience}

	\textbf{Resident Researcher} \hfill 2025-Present \\
	Department of Radiation Medicine and Applied Sciences, UC San Diego
	\vspace{-.2cm}
	\begin{itemize}[label=$\cdot$]
		\itemsep-0.5em
		\item Managed 5 concurrent clinic-forward research projects spanning MLC-based radiosurgery technique development, treatment planning system and algorithm commissioning, and clinical workflow automation, prioritized to address emerging clinical needs \textit{(see Clinical Projects)}.
		\item Prepared and delivered educational materials for physician and physics faculty audiences.
	\end{itemize}

	\textbf{Doctoral Researcher} \hfill 2020-2025 \\
	Department of Physics \& Astronomy, Louisiana State University
	\vspace{-.2cm}
	\begin{itemize}[label=$\cdot$]
		\itemsep-0.5em
		\item Designed a prototype cardiovascular digital twin capable of modeling up to 34 billion vessels across the full human body using distributed supercomputing resources.
		\item Developed parallel-enabled biophysics models of blood flow, vessel deformation and rupture, and vasodilation using MPI and Kokkos runtime systems for exascale computing.
		\item Performed computational feasibility analyses of simulating radiation-induced vascular injury and hemodynamic changes at full-body scale, informing subsequent model development.
		\item Served as principal point of contact for a multi-institutional research collaboration spanning Los Alamos National Laboratory, Oak Ridge National Laboratory, and the RIKEN Center for Computational Science.
		\item Wrote and contributed to grant proposals, including two NSF proposals totaling over \$1.8 million in awarded funding (PI: Hartmut Kaiser).
		\item Mentored 3 undergraduate and 2 master's students in research methods, data analysis, and scientific communication.
	\end{itemize}

\newpage
	\textbf{Undergraduate Research Assistant} \hfill 2019-2020 \\
	Department of Physics \& Astronomy, Louisiana State University
	\vspace{-.2cm}
	\begin{itemize}[label=$\cdot$]
		\itemsep-0.5em
		\item Developed personalized anthropomorphic phantoms for radiation dose measurement and quality assurance.
		\item Organized research meetings and facilitated cross-institutional research partnerships.
		\item Conducted literature reviews to support ongoing projects and inform research directions.
	\end{itemize}

\end{rSection}
\vspace{0.5cm}

%----------------------------------------------------------------------------------------
%	Publications and Presentations
%----------------------------------------------------------------------------------------

\begin{rSection}{Publications \& Presentations}

	\textbf{Manuscripts Submitted for Publication}
	\begin{enumerate}
		\itemsep-0.3em
		\item \textbf{Cole, M}, Newhauser W, Diehl P, and Seleson P. ``Computational feasibility and scalability of coupling peridynamic vessel-fracture simulation with a whole-body cardiovascular digital twin.'' Submitted to \textit{Biomechanics and Modeling in Mechanobiology}, August 2026.
		\item \textbf{Cole, M}, Newhauser W, and Diehl P. ``A multiscale biophysical model of nitric oxide diffusion and vasodilation: A computational feasibility study.'' Submitted to \textit{Biomechanics and Modeling in Mechanobiology}, August 2026.
	\end{enumerate}

	\textbf{Preprints}
	\begin{enumerate}
		\itemsep-0.3em
		\item Newhauser W, \textbf{Cole, M}, Kaiser H, Tohid R, Chancellor J, Nader N, and Diehl P. ``Toward Large-Scale Numerical Modeling of the Cardiovascular System with up to 34 Billion Vessels.'' \textit{bioRxiv} 2026.05.22.727287 (2026).
	\end{enumerate}

	\textit{3 additional manuscripts in preparation; titles and details available upon request.}

	\textbf{Invited Talks \& Conference Presentations}
	\begin{enumerate}
		\itemsep-0.3em
		\item \textbf{Cole, M}. \textit{First Digital Biophysical Model of the Entire Human Cardiovascular System} [Invited Talk]. Advances in Applied Computer Science Invited Speaker Series, Los Alamos National Laboratory, Los Alamos, NM. August 27, 2025.
		\item \textbf{Cole, M}. \textit{Simulating Space Radiation-Induced Changes to Cardiovascular Blood Flow} [Poster]. Aerospace Medical Association Annual Scientific Meeting, New Orleans, LA. May 25, 2023.
		\item \textbf{Cole, M}. \textit{Computational Feasibility of Simulating Radiation-Induced Changes in Vasculature and Blood Flow Rates in the Entire Human Body} [Poster]. Workshop on Asynchronous Many Task Applications, Baton Rouge, LA. February 15, 2023.
	\end{enumerate}

\end{rSection}
\vspace{0.5cm}

%----------------------------------------------------------------------------------------
%	TEACHING EXPERIENCE
%----------------------------------------------------------------------------------------

\begin{rSection}{Teaching Experience}

	\textbf{Joint Conference Case Lecture} \hfill 2026 \\
	UC San Diego
	\vspace{-.2cm}
	\begin{itemize}[label=$\cdot$]
		\itemsep-0.5em
		\item Presented a lecture on the radiation treatment of trigeminal neuralgia.
		\item Collaborated with a radiation oncology resident to co-develop and deliver the lecture to clinical faculty.
	\end{itemize}

	\textbf{Resident Lecture Series} \hfill 2025-2026 \\
	UC San Diego
	\vspace{-.2cm}
	\begin{itemize}[label=$\cdot$]
		\itemsep-0.5em
		\item Presented a series of lectures to physics faculty covering:
		\vspace{-.2cm}
		\begin{itemize}[label=$\circ$]
			\itemsep-0.5em
			\item Proton QA
			\item MR-Only Simulation and Planning
			\item Commissioning and QA of Treatment Planning Systems
			\item TG-51
			\item Biennial QA for Halcyon and Ethos
			\item Treatment Planning Considerations for Gastrointestinal SBRT
			\item TG-29: Total Body Irradiation
		\end{itemize}
	\end{itemize}

	\textbf{Guest Lecturer} \hfill 2022 \\
	Louisiana State University
	\vspace{-.2cm}
	\begin{itemize}[label=$\cdot$]
		\itemsep-0.5em
		\item Presented a 5-lecture course module on research methods and technologies to graduate students.
		\item Designed the module curriculum and delivered real-time coding demonstrations.
		\item Prepared in-class and take-home assignments to assess student comprehension.
	\end{itemize}

	\textbf{Teaching Assistant / Learner's Assistant} \hfill 2019-2021 \\
	Louisiana State University
	\vspace{-.2cm}
	\begin{itemize}[label=$\cdot$]
		\itemsep-0.5em
		\item Delivered lectures, held office hours, and graded assignments for undergraduate physics courses (2020-2021).
		\item Tutored underclassmen and led review sessions to reinforce course material (2019-2020).
	\end{itemize}

\end{rSection}
\vspace{0.5cm}

%----------------------------------------------------------------------------------------
%	Awards, Funding, and Other Support
%----------------------------------------------------------------------------------------

\begin{rSection}{Honors \& Awards}

	\textbf{Dr. Charles M. Smith Superior Graduate Student Scholarship}; \$20,000 \hfill 2024 \\
	Louisiana State University

	\textbf{Roussel Family Graduate Student Award in Communication}; \$2,000 \hfill 2023 \\
	Louisiana State University

	\textbf{Best Poster Award} \hfill 2023 \\
	Workshop on Asynchronous Many Task Applications

	\textbf{Kenneth R. Hogstrom Superior Graduate Student Scholarship Award}; \$20,000 \hfill 2022 \\
	Louisiana State University

	\textbf{LaSPACE Graduate Student Research Award}; \$16,000 \hfill 2022 \\
	Louisiana Space Grant Consortium

	\textbf{Inspire Award for Exemplary Ethical Action}; \$1,000 \hfill 2020 \\
	LSU Ethics Institute

\end{rSection}
\vspace{0.5cm}


%----------------------------------------------------------------------------------------
%	Outreach and Volunteer Work
%----------------------------------------------------------------------------------------

\begin{rSection}{Outreach and Volunteer Work}

	\textbf{Project Co-Lead}, Medical Physics Brain (Rayos Contra Cancer) \hfill 2026-Present
	\vspace{-.2cm}
	\begin{itemize}[label=$\cdot$]
		\itemsep-0.5em
		\item Lead and manage an effort to build an open-access medical physics education repository.
		\item Consult with and recruit practicing medical physicists to contribute clinical and technical content.
		\item Design and architect the platform's structure and content organization.
		%\item Author and edit reference pages for an open-access repository of medical physics educational material.
	\end{itemize}

%	\textbf{Assistant Coach}, Special Olympics of Southern California \hfill 2025-Present
%	\vspace{-.2cm}
%	\begin{itemize}[label=$\cdot$]
%		\itemsep-0.5em
%		\item [FILL IN LATER: sport/team, coaching duties, frequency]
%	\end{itemize}

	\textbf{Advocate/Judge}, Louisiana Science \& Engineering Fair, Region VII \hfill 2022
	\vspace{-.2cm}
	\begin{itemize}[label=$\cdot$]
		\itemsep-0.5em
		\item Judged science projects for middle and high school students.
		\item Served as a scientific advocate, providing students with resources and feedback to advance their projects.
	\end{itemize}

	\textbf{Co-Founder}, One Louisiana NOW! \hfill 2020
	\vspace{-.2cm}
	\begin{itemize}[label=$\cdot$]
		\itemsep-0.5em
		\item Co-founded a non-profit relief organization to supply personal protective equipment to hospitals in need during the COVID-19 pandemic.
		\item Researched and designed prototype protective gear based on hospital-identified needs.
		\item Coordinated with institutions and hospital staff to ensure timely delivery of supplies.
		\item Co-authored an unpublished internal report, ``A Community-Based Response to Equipment Shortages Caused by the COVID-19 Pandemic,'' circulated among participating hospitals and physicians; the report was subsequently used by other organizations as a model to establish their own PPE response programs.
	\end{itemize}

\end{rSection}
\vspace{0.5cm}


%----------------------------------------------------------------------------------------
%	ASSOCIATIONS / SERVICE
%----------------------------------------------------------------------------------------

\begin{rSection}{Service}

	\textbf{Professional Organizations}
	\vspace{-.2cm}
	\begin{itemize}[label=$\circ$]
		\itemsep-0.5em
		\item \textbf{American Association of Physicists in Medicine (AAPM)} \hfill 2019-Present
		\item \textbf{Radiation Research Society (RRS)} \hfill 2022-Present
	\end{itemize}

	\textbf{Journal Peer Reviewer}
	\vspace{-.2cm}
	\begin{itemize}[label=$\circ$]
		\itemsep-0.5em
		\item \textit{Medical Physics}, AAPM \hfill 2026
		\item \textit{PLOS One} \hfill 2025
	\end{itemize}

	\textbf{Institutional Committees}
	\vspace{-.2cm}
	\begin{itemize}[label=$\circ$]
		\itemsep-0.5em
		\item Quality Committee, UC San Diego \hfill 2026-Present
		\item Communications Committee, UC San Diego \hfill 2026-Present
	\end{itemize}

\end{rSection}
\vspace{0.5cm}


%----------------------------------------------------------------------------------------
%	TECHNICAL SKILLS
%----------------------------------------------------------------------------------------

\begin{rSection}{Technical Skills}

	\begin{description}
		\itemsep-0.4em
		\item[Treatment Planning \& Oncology Information Systems:] Eclipse, Ethos, ARIA, ESAPI
		\item[Programming Languages:] Python, C, C\#, C++, MATLAB, Swift
		\item[Software \& Systems Development:] iOS app development (Swift/SwiftUI), web development, Git/version control, embedded systems (Raspberry Pi)
		\item[HPC Systems:] scripting, job management systems, parallel computing (\textit{e.g.}, MPI, Kokkos)
		\item[Data Visualization:] matplotlib, gnuplot, ParaView
	\end{description}

\end{rSection}
%\vspace{0.5cm}

%----------------------------------------------------------------------------------------
%	REFERENCES
%----------------------------------------------------------------------------------------

%\begin{rSection}{References}
%Available upon request.
%\end{rSection}

\end{document}
